\documentclass{article}
\usepackage[utf8]{inputenc}

\usepackage{tikz}
\usepackage{tikz-cd}
\usepackage{graphicx} % Required for inserting images
\usepackage[a4paper, margin=1in]{geometry}
\usepackage{amsfonts}
\usepackage{amsmath} % Set custom margins
\usepackage{parskip}
\usepackage{amssymb} % disable indentation
\usepackage{amsthm}
\usepackage{annotate-equations}


\title{Category Theory - Lecture 6 (Notes)}
\author{Vorashil Farzaliyev}
\date{September 2024}

% Define a new theorem style that starts content on a new line


\newtheorem{remark}{Remark}[section]
\newtheorem{lemma}{Lemma}[section]
\newtheorem{proposition}{Proposition}[section]
\newtheorem{definition}{Definition}[section]
\newtheorem{observation}{Observation}[section]
\newtheorem{example}{Example}[section]
\newtheorem{exercise}{Exercise}[section]
\newtheorem{theorem}{Theorem}[section]
\newtheorem{note}{Note}[section]

\renewcommand{\shorto}{\to}
% new command for making \mapsto 0.5cm

\renewcommand{\to}{\xrightarrow{\hspace{0.5cm}}}  % Adjust the length as needed
\renewcommand{\mapsto}{\mathrel{\kernel\longmapsto\kernel}}  % Adjust spacing with \mkern if needed

\begin{document}
    \maketitle

    From last time...


    \begin{proposition}
        Let \(\phi: F \Rightarrow G\) be a natural transformation.

        \[
            \begin{tikzcd}[column sep=7em, row sep=4em]
                \mathbf{C} \arrow[r, bend left=50, "F"{name=F, above}]
                \arrow[r, bend right=50, "G"{name=G, below}]
                & \mathbf{D}
                \arrow[Rightarrow, from=F, to=G, "\phi"]
            \end{tikzcd}
        \]

        Assume that for all \(A \in \mathbf{C}\), \(\phi_A: FA \to GA\) is an isomorphism with
        inverse \(\phi^{-1}_A: GA \to FA\).

        Then \(\phi^{-1}: G \Rightarrow F\) is defined by

        \[
            (\phi^{-1})_A = (\phi_A)^{-1} \text{ for all } A \in \mathbf{C}
        \]
    \end{proposition}

    \begin{proof}
        We show that the naturality, as a family of maps, is given by

        \[
            \begin{array}{c c} % Create a 2x2 array
                \mathbf{C} & \mathbf{D} \\
                \begin{tikzcd}[column sep=4em, row sep=4em]
                    A \arrow[d, mapsto, "f"] \\
                    B
                \end{tikzcd}
                &
                \begin{tikzcd}[column sep=7em, row sep=4em]
                    GA \arrow[r, "\phi^{-1}_A"] \arrow[d, "Gf"]
                    & FA \arrow[d, "Ff"] \\
                    GB\arrow[r, "\phi^{-1}_B"]
                    & FB
                \end{tikzcd}
                &
                \begin{tikzpicture}[overlay]
                    % Arrow to the first rectangle (FA, GA, FB, GB)
                    \draw[->, thick] (3, -1) -- (3, 1.5) node[near start, auto] {But does this commute?};
                    r
                \end{tikzpicture}
            \end{array}
        \]

        \vspace{0.2in}
        We have

        \[
            \begin{tikzcd}[column sep=5em, row sep=4em]
                GA \arrow[r, "\phi^{-1}_A"] \arrow[d, "Gf"]
                \arrow[rr, bend left=35, "1_{GA}"] % Curved arrow from FA to HA
                & FA \arrow[d, "Ff"] \arrow[r, "\phi_A"]
                & GA \arrow[d, "Gf"] \\
                GB \arrow[r, "\phi^{-1}_B"]
                \arrow[rr, bend right=35, "1_{GB}"] % Curved arrow from FB to HB
                & FB \arrow[r, "\phi_B"]
                & GB
            \end{tikzcd}
            &
            \begin{tikzpicture}[overlay]
                % Arrow to the first rectangle (FA, GA, FB, GB)
                \node at (-4.5, 0) {(1)};

            \end{tikzpicture}
        \]

        \newpage

        We show that \((1)\) commutes by post-composing it with
        \(\phi_B\), \(\phi_A\), \(1_{GA}\), \(1_{GB}\) and
        showing the following equality holds.

        \[
            \phi_A^{-1} Ff = Gf \phi_B^{-1}
        \]

    \end{proof}


    \section{Equivalence of categories}

    \subsection{Motivation}

    We often have

    \[
        \begin{tikzcd}[column sep=7em, row sep=4em]
            \mathbf{C} \arrow[r, bend left=50, "F"{name=F, above}]
            & \mathbf{D} \arrow[l, bend left=50, "G"{name=G, above}]
        \end{tikzcd}
    \]

    but asking
    \begin{gather*}
        GFA = A \text{ for all } A \in \mathbf{C}\\
        FGX = X \text{ for all } X \in \mathbf{D}\\
    \end{gather*}

    is too restrictive. Because we are interested in isomorphisms, not the equality of objects


    \begin{definition}
        Let \(F: \mathbf{C} \to \mathbf{D}\) be a functor.

        We say \(F\) is an \underline{equivalence} if there exists a functor \(G: \mathbf{D} \to \mathbf{C}\) and
        following natural transformations
        \begin{gather*}
            \eta: 1_{\mathbf{C}} \Rightarrow GF\\
            \xi: FG \Rightarrow 1_{\mathbf{D}}\\
        \end{gather*}

    \end{definition}

    \begin{note}
        The components of \(\eta\) and \(\xi\) are isomorphisms.

        \[
            \begin{array}{c c} % Create a 2x2 array
                \underbrace{A \stackrel{\eta_A}{\xrightarrow{\hspace{1cm}}} GFA}_{in \ \mathbf{C}} \quad
                &
                \underbrace{FGX \stackrel{\xi_X}{\xrightarrow{\hspace{1cm}}} X}_{in \ \mathbf{D}} \\

                \forall A \in \mathbf{C} & \forall X \in \mathbf{D} \\
            \end{array}
        \]
    \end{note}

    \vspace{0.2in}

    \begin{theorem}
        Let \(F: \mathbf{C} \to \mathbf{D}\) be a functor.

        Then \(F\) is an equivalence if and only if \(F\) is essentially surjective and fully faithful.
    \end{theorem}

    \begin{proof}
        \((\Rightarrow)\): Exercise.

        \((\Leftarrow)\): Assume \(F\) is essentially surjective and fully faithful.

        \[
            \begin{align}
                \text{Assume \(F\) is essentially surjective} &\Rightarrow (\forall X \in \mathbf{D})  (\exists A \in \mathbf{C})  (\exists \text{ isomorphism } FA \to X) \\
                & \quad \ \quad \ \quad \  \downarrow \text{(axiom of choice)}  \\
                &\Rightarrow \text{we have a function } G: Ob(\mathbf{D}) \to Ob(\mathbf{C}) \\
                & \quad \quad \text{ a family of maps } (\xi_X: FGX \to X | X \in \mathbf{D}) \\
                & \quad \quad \text{ such that } \xi_X: FGX \to X \text{ is an isomorphism } \forall X \in \mathbf{D} \\
            \end{align}
        \]

        \textbf{Task 1:} Extend \(G\) to a functor \(G\) and \(\xi\) to natural transformation.

        For \(f: X \to Y\) in \(\mathbf{D}\), we need to define \(G(f): GX \to GY\) in \(\mathbf{C}\).
        We use the fact that \(F\) is fully faithful

        \[
            \mathbf{C}(GX, GY) \stackrel{F}{\to} \mathbf{D}(FGX, FGY) \text{ is bijection.}
        \]

        This means for every \(FGX \stackrel{v}{\to} FGY\) there is a unique \(GX \stackrel{u}{\to} GY\) s.t
        \[
            FGX \stackrel{Fu}{\to} FGY = FGX \stackrel{v}{\to} FGY
        \]

        So to show that \(G(gf) = G(g) G(f)\)), we show \(G(g) G(f)\) has the same property, i.e

        \[
            \begin{align}
                F(G(g) G(f)) : FGX \to FGZ = FGX \stackrel{\xi_X}{\to} X \stackrel{gf}{\to} Z \stackrel{\xi_Z^{-1}}{\to} FGZ \\
            \end{align}
        \]

        We know
        \[
            \begin{tikzcd}[column sep=7em, row sep=4em]
                FGX \arrow[r,  "\xi_X"] \arrow[d, "FG(gf)"]
                & X \arrow[d, "gf"] \\
                FGZ \arrow[r, "\xi_Z"]
                & Z
            \end{tikzcd}
        \]

        Consider
        \[
            FGX \stackrel{\xi_X}{\to} X \stackrel{f}{\to} Y \stackrel{\xi_Y^{-1}}{\to} FGY
        \]

        So there is a unique map, which we write \(Gf: GX \to GY\) such that

        \[
            FGf = FGX \to FGY = FGX \stackrel{\xi_X}{\to} X \stackrel{f}{\to} Y \stackrel{\xi_Y^{-1}}{\to} FGY \ \text{ (*) }
        \]

        Here

        \[
            \begin{array}{c c} % Create a 2x2 array
            (*)
                \Leftrightarrow &
                \begin{tikzcd}[column sep=7em, row sep=4em]
                    FGX \arrow[r,  "\xi_X"] \arrow[d, "FG(gf)"]
                    & X \arrow[d, "gf"] \\
                    FGZ \arrow[r, "\xi_Z"]
                    & Z
                \end{tikzcd} \text{ commutes } \\

            \end{array}
        \]

        To check \(G\) is a functor, let \(f: X \to Y\) and \(g: Y \to Z\) be in \(\mathbf{D}\) and show

        \[
            G(gf): GX \to GZ = G(g) G(f): GX \to GZ
        \]

        By definition, \(G(gf)\) is the unique map

        \[
            u: GX \to GZ
        \]
        such that
        \[
            F(u): FGX \to FGZ = FGX \stackrel{\xi_X}{\to} X \stackrel{f}{\to} Z \stackrel{\xi_Z^{-1}}{\to} FGZ
        \]

        We need to show following diagram commutes

        \[
            \begin{tikzcd}[column sep=7em, row sep=4em]
                FGX \arrow[r,  "\xi_X"] \arrow[d, "F(Gg Gf)"]
                & X \arrow[d, "gf"] \\
                FGZ \arrow[r, "\xi_Z"]
                & Z
            \end{tikzcd}
        \]
        But this follows from
        \[
            \begin{tikzcd}[column sep=7em, row sep=4em]
                FGX \arrow[r, "\xi_X"] \arrow[d, "FGf"]
                & X \arrow[d, "f"] \\
                FGY \arrow[r, "\xi_Y"] \arrow[d, "FGg"]
                & Y \arrow[d, "g"] \\
                FGZ \arrow[r, "\xi_Z"]
                & Z
            \end{tikzcd}
            \begin{tikzpicture}[overlay]
                % Arrow to the first rectangle (FA, GA, FB, GB)
                \node at (-2, 1) {(1)};
                \node at (-2, -1) {(2)};

            \end{tikzpicture}
        \]

        Since both \((1), (2)\) commute.

        In the same way we show

        \[
            G(1_X) = 1_{GX}
        \]
    \end{proof}
\end{document}
